% Part: first-order-logic
% Chapter: beyond
% Section: introduction

\documentclass[../../../include/open-logic-section]{subfiles}

\begin{document}

\olfileid[tr]{fol}{byd}{int}

\olsection{Genel Bakış}

Birinci dereceden mantık, ilgi çekici tek mantık sistemi değildir: onun birçok
genişletmesi ve çeşitlemesi vardır. Bir mantık tipik olarak bir dilin biçimsel
belirleniminden, çoğu zaman fakat her zaman değil bir türetim sisteminden ve
yine çoğu zaman fakat her zaman değil amaçlanan bir semantikten oluşur. Ancak
``mantık'' teriminin bu teknik kullanımı açık bir soru doğurur: birinci
dereceden mantık olmayan mantıkların, sözcüğün sezgisel ya da felsefi
anlamındaki ``mantık''la ilişkisi nedir? Aşağıda betimlenen sistemlerin tümü
şu ya da bu türden akıl yürütmeyi modellemek üzere tasarlanmıştır; onları
mantıksal kılan şeyin ne olduğunu söyleyebilir miyiz?

Bu soruların kolay yanıtları yoktur. ``Mantık'' sözcüğü farklı biçimlerde ve
farklı bağlamlarda kullanılır; kavram, ``doğruluk'' kavramı gibi, çok sayıda
felsefi konumdan çözümlenmiştir. Örneğin mantıksal akıl yürütmenin amacının,
hangi ifadelerin zorunlu olarak doğru, a priori doğru, mantık dışı terimlerin
yorumundan bağımsız olarak doğru, biçimleri gereği doğru ya da dilsel uzlaşım
gereği doğru olduğunu belirlemek olduğu düşünülebilir; bu anlayışların her
biri ciddi ölçüde açıklığa kavuşturulmayı gerektirir. Dikkatimizi matematikte
kullanılan mantık türüyle sınırlasak bile kapsamı konusunda pek az görüş
birliği vardır. Örneğin Russell ile Whitehead, \textit{Principia Mathematica}
içinde, Frege'nin başlattığı {\em mantıkçı} gelenek doğrultusunda matematiği
mantık temelinde geliştirmeye çalıştılar. Kullandıkları mantık sistemi,
aşağıda betimlenene benzeyen bir yüksek tipli mantık biçimiydi. Sonunda çoğu
ölçüte göre salt mantıksal görünmeyen aksiyomlar (özellikle sonsuzluk aksiyomu
ile indirgenebilirlik aksiyomu) eklemek zorunda kaldılar; yine de bazı yüksek
dereceden akıl yürütme biçimlerinin mantıksal sayılması gerektiği savunulabilir.
Buna karşılık ontolojisinde ``önermeleri'' meşru söylem nesneleri olarak kabul
etmeyen Quine, ikinci ve yüksek dereceden mantığın aslında koyun postuna
bürünmüş küme kuramından ibaret olduğunu savunur; başka deyişle yüklemler
üzerinde niceleme içeren sistemler salt mantıksal değildir.

Şimdilik bu tür felsefi sorunları başka bir güne bırakmak ve aşağıdaki
sistemleri, mantıksal olsun ya da olmasın, çeşitli akıl yürütme türlerinin
biçimsel idealleştirmeleri olarak düşünmek en iyisidir.

\end{document}
